% ===== Appendix: Corrections Log =====
\clearpage\appendix
\section{Corrections Log}\label{sec:corrections}
The original handwritten notes (which the syllabus itself warns may contain typos) were corrected in the places listed below.
Each entry gives the page of the original notes (the gray \textit{p.\,N} margin tag), the section of these notes where the correction appears, what the handwritten notes had, and the correction together with a brief reason.
Purely cosmetic changes (layout, re-ordering, added explanations) are not listed.

{\small
\setlength{\tabcolsep}{4pt}
\renewcommand{\arraystretch}{1.2}
\newcommand{\grp}[1]{\multicolumn{4}{@{}l}{\rule{0pt}{1.4em}\textbf{\S\ref{#1}\enspace\nameref{#1}}}\\*[2pt]}
\begin{longtable}{@{}>{\centering\arraybackslash}p{\dimexpr(\linewidth-6\tabcolsep)*62/1000\relax}>{\centering\arraybackslash}p{\dimexpr(\linewidth-6\tabcolsep)*58/1000\relax}>{\raggedright\arraybackslash}p{\dimexpr(\linewidth-6\tabcolsep)*347/1000\relax}>{\raggedright\arraybackslash}p{\dimexpr(\linewidth-6\tabcolsep)*533/1000\relax}@{}}
\toprule
\textbf{Orig.\ page} & \textbf{Sec.} & \textbf{Notes had} & \textbf{Corrected to / reason}\\
\midrule
\endfirsthead
\multicolumn{4}{@{}l}{\footnotesize\itshape Corrections Log (continued)}\\
\toprule
\textbf{Orig.\ page} & \textbf{Sec.} & \textbf{Notes had} & \textbf{Corrected to / reason}\\
\midrule
\endhead
\midrule
\multicolumn{4}{r@{}}{\footnotesize\itshape continued on next page}\\
\endfoot
\bottomrule
\endlastfoot
%
\grp{sec:fields-vs}
2 & \ref{ss:fields} & $x\oplus y = x_1+y_1+j(y_1+y_2)$ & $x\oplus y = x_1+y_1+j(x_2+y_2)$: the second component is $x_2+y_2$.\\
3 & \ref{ss:vector-spaces} & ``$\Rightarrow$'' between the members of the chain proving $y\oplus y=y$ & ``$=$'': the steps are equalities.\\
5 & \ref{ss:subspaces} & Middle graph labelled ``$\in\W$'' & $\in\W_2$ (the set of functions in the example).\\
6 & \ref{ss:sums} & ``$Y\cup Z$ is not closed under addition but not closed under scalar multiplication'' & $Y\cup Z$ is closed under scalar multiplication but (in general) not under addition.\\
6 & \ref{ss:sums} & $w_1\oplus_{\V} w_2$ and $\alpha\odot_{\V} w_1$ in $\V\times\W$ & $w_1\oplus_{\W} w_2$, $\alpha\odot_{\W} w_1$: the second components use the operations of $\W$.\\
9 & \ref{ss:basis} & ``some of the vectors are nonzero'' & Some of the \emph{coefficients} $\alpha_i-\beta_i$ are nonzero.\\
10 & \ref{ss:change-basis} & ``$P=[w_1]_{B_1}$'', ``$P=[w_2]_{B_1}$'' and ``$P\cdot[w_3]_{B_2}$'' on the right & Each is the corresponding column of $P$, i.e.\ $[w_j]_{B_1}$.\\
\grp{sec:normed}
11 & \ref{ss:norms} & $\norm{x_i}$ in step (1); sum without parentheses in step (3) & $\abs{x_i}$; $\sum_i\big(\abs{x_i}+\abs{y_i}\big)$.\\
11 & \ref{ss:lp} & 1st quadrant: $x_1-x_2=1$ & $x_1+x_2=1$, since $\abs{x_1}+\abs{x_2}=x_1+x_2$ there; 3rd/4th quadrant cases added.\\
12 & \ref{ss:function-norms} & $\mathcal{L}_p$ and $\mathcal{L}_2$ integrals without $dt$; $\mathcal{L}_1$ norm without subscript; $\max_t\abs{f(t)}$ for $\mathcal{L}_\infty$ & $dt$ and subscript added; $\sup_t\abs{f(t)}$ (the example below shows why the supremum is needed).\\
%
13 & \ref{ss:matrix-norms} & $\tr(A)$ on the left-hand side & $\tr(A\T A)$: the trace taken is that of $A\T A$.\\
15 & \ref{ss:matrix-norms} & Stray ``7'' written below the ``8'' & $\norm{A}_1=8$ (column sums $6$ and $8$); $7$ is $\norm{A}_\infty$ from the previous page.\\
15 & \ref{ss:convergence} & $\norm{(-1)^n-1}=(-1)^n-1$ & $\abs{(-1)^n-1}$; without the absolute value the expression is negative for odd $n$.\\
15 & \ref{ss:convergence} & $\mathcal{L}_1$ integral over $(-\infty,\infty)$ & $\int_{-1}^{1}$: functions in $\V$ live on $[-1,1]$.\\
16 & \ref{ss:cauchy} & Inequality step justified ``by absolute convergence'' & By the triangle inequality (and $\norm{v-v_m}=\norm{v_m-v}$ by homogeneity with $\alpha=-1$).\\
18 & \ref{ss:cauchy} & $\lim_{n,m}\norm{f_n-f_m}_\infty=\abs{1-\frac nm}$ & The distance itself is $\norm{f_n-f_m}_\infty=\abs{1-\frac nm}$ ($m\ge n$), which does not tend to $0$.\\
%
\grp{sec:ips}
19 & \ref{ss:inner-products} & Row vector $[y_1\ y_2\ \cdots]$ in $y\Hm x$ & $[\bar y_1\ \bar y_2\ \cdots]$: $y\Hm=\bar y\T$ has conjugated entries.\\
19 & \ref{ss:inner-products} & $B\T$ with rows $b_{11}\,b_{12}$ / $b_{12}\,b_{22}$ / $b_{13}\,b_{23}$ & Entry $(1,2)$ of $B\T$ is $b_{21}$.\\
19 & \ref{ss:cauchy-schwarz} & Equality iff $x=cy$ & Also $y=0_\V$ ($x$ arbitrary) gives $0=0$; equality iff $x,y$ are linearly dependent.\\
24 & \ref{ss:gram-schmidt} & $w_2=\bar w_2/\norm{w_2}$ & $w_2=\bar w_2/\norm{\bar w_2}$; the last step $w_n=\bar w_n/\norm{\bar w_n}$ (implied by the pattern) added.\\
%
\grp{sec:lt}
25 & \ref{ss:lt-def} & Sums for $v_1,v_2$ starting at $i=1$ & $\sum_{i=0}^{n}$: the constant term must be included (the next line already uses $i=0$).\\
26 & \ref{ss:null} & $\Null(\mathcal{A}_1)=0_\V$ & $\Null(\mathcal{A}_1)=\{0_\V\}$: the null space is a set.\\
28 & \ref{ss:one-to-one} & Side note labelled $\Null(\mathcal{A})=\{\text{poly.\ degree}\le1\}$ & This is $\Range(\frac{d}{dt})$ (cf.\ p.~29); the null space is the constants.\\
28 & \ref{ss:range} & $\alpha w_1=\mathcal{A}(\alpha v_1)=\alpha\mathcal{A}v_1=\alpha w_1\in\Range(\mathcal{A})$ (circular chain) & Rewritten as the implication $\mathcal{A}(\alpha v_1)=\alpha w_1\Rightarrow\alpha w_1\in\Range(\mathcal{A})$.\\
29 & \ref{ss:range} & $\dim\Range(\mathcal{A}_1)=0$ in the last line & $\dim\Range(\mathcal{A}_2)=0$.\\
29 & \ref{ss:range} & $\Range(\mathcal{A})\not\subseteq\W$ (definition of ``into'') & $\Range(\mathcal{A})\subsetneq\W$: $\Range(\mathcal{A})$ is always a subset of $\W$, a proper subset is meant.\\
31 & \ref{ss:basis-change-rep} & $[v]_B$ (no bar) in the second line & $[v]_{\bar B}$; the last two lines (combining 1) and 2), giving $\bar A=Q^{-1}AP$) added.\\
32 & \ref{ss:basis-change-rep} & First line $\bar A=Q^{-1}\bar A$ & $\bar{\bar A}=Q^{-1}\bar A$ (double bar on the left).\\
%
\grp{sec:direct-sum}
34 & \ref{ss:orth-complement} & ``$M_2$ \& $M_2$'' & $M_1$ \& $M_2$.\\
35 & \ref{ss:orth-complement} & Un-orthogonalized $[1\ 0\ 0\ 0]\T$, $[0\ 1\ 0\ 0]\T$ labelled $e_3,e_4$, and $M^\perp=\Span\{e_3,e_4\}$ & These are not orthogonal to $M$ (e.g.\ $\inner{[1\ 0\ 0\ 0]\T}{m_1}=1$); Gram--Schmidt computation added (verified in exact arithmetic). The cramped last vector is read as $[0\ 1\ 0\ 0]\T$.\\
40 & \ref{ss:gram} & Third basis matrix labelled $v_2$ & $v_3$.\\
41 & \ref{ss:proj-matrices} & $Bx\Hm=[m_i\Hm x]=[\inner{x}{m_i}]=[\text{Gram matrix}]$ & $B\Hm x=[\inner{x}{m_i}]=b$; the Gram matrix is $B\Hm B$ (entry $(i,j)=\inner{m_j}{m_i}$), moved to the line above.\\
41 & \ref{ss:proj-matrices} & $P_{M^\perp}=(I-P_M)x$ & $P_{M^\perp}=I-P_M$; $(I-P_M)x$ is the projected vector.\\
41 & \ref{ss:proj-matrices} & $B(B\Hm B^{-1})B\Hm$ in the last step & $B(B\Hm B)^{-1}B\Hm$.\\
%
\grp{sec:axb}
42 & \ref{ss:exist-unique} & Second small ``$\C^m$'' label next to the zero vector in the figure & Read as the subscript $0_{\C^m}$, consistent with $\Null(A)=\{x: Ax=0_{\C^m}\}$.\\
43 & \ref{ss:exist-unique} & ``$=b+0=0$'' at the end; ``$x$ is exact'' & $=b+0=b$; $\hat x$ is exact.\\
44 & \ref{ss:least-squares} & $\norm{Ax-b_1}+\norm{b_2}^2$ & $\norm{Ax-b_1}^2+\norm{b_2}^2$ (square missing).\\
45 & \ref{ss:least-squares} & Projection applied to $[1\ 1]\T$ & Applied to $b=[2\ 1]\T$, giving $b_1=[3/2\ 3/2]\T$.\\
49 & \ref{ss:rank-cases} & $b\in\C^m\ \forall b\in\Range(A)$ & $b\in\Range(A)\ \forall b\in\C^m$ (quantifiers swapped).\\
49 & \ref{ss:rank-cases} & $\operatorname{Proj}_{\Range(A\Hm)}\hat x=\operatorname{Proj}_{\Range(A\Hm)}\hat x$ & $\operatorname{Proj}_{\Null^\perp(A)}\hat x=\operatorname{Proj}_{\Range(A\Hm)}\hat x$.\\
50 & \ref{ss:rank-cases} & $\hat x+m$ & $\hat x+n$, $n\in\Null(A)$ (an element of the null space is added).\\
50 & \ref{ss:rank-cases} & $z=[x_2\ x_3]\T$ & $z=[x_1\ x_2]\T$: $A_1$ consists of the first two columns of $A$ (consistent with the final $\hat x$).\\
%
\grp{sec:spectral}
51 & \ref{ss:invariant} & ``$\notin M$'' & $\notin M_1$ (made explicit).\\
53 & \ref{ss:eigen} & $\tr\mathcal{A}\triangleq\det A$ & $\tr\mathcal{A}\triangleq\tr A$: the trace of the operator is the trace of its matrix.\\
53 & \ref{ss:eigen} & ``$R$: rotation vectors'' & $R$ denotes the rotation matrix.\\
56 & \ref{ss:block-diag} & $M_1=\Span\{b_1^1,b_2^2\}$, $M_2=\Span\{b_1^1,b_2^2,b_3^3\}$ & $M_1=\Span\{b_1^1,b_1^2\}$, $M_2=\Span\{b_2^1,b_2^2,b_2^3\}$ ($b_i^j$: $j$-th basis vector of $M_i$), matching $P_1,P_2$.\\
56 & \ref{ss:block-diag} & $M_2=\{b_2^1\}$ & $M_2=\Span\{b_2^1\}$ (span, not the set).\\
57 & \ref{ss:block-diag} & $P=[b_1^1\ b_2^1\mid b_2^1]$ & $P=[b_1^1\ b_1^2\mid b_2^1]$: basis of $M_1$ followed by that of $M_2$, as in the next line.\\
57 & \ref{ss:block-diag} & Small mark ($*$, $+$) on the circled exponent 2 & Dropped (it only linked the exponent to the multiplicity in $d(s)$).\\
57 & \ref{ss:diagonalization} & $(A-\lambda_2I)e_2$ in the second term (which would be $0$) & $(A-\lambda_1I)e_2=(\lambda_2-\lambda_1)e_2\neq0$.\\
57 & \ref{ss:diagonalization} & $A\in\C^n$ & $A\in\C^{n\times n}$.\\
58 & \ref{ss:diagonalization} & Diagonal matrix labelled $A$ & $\bar A=P^{-1}AP$.\\
%
\grp{sec:polys}
59 & \ref{ss:charpoly} & $A_{31}=1$ & $A_{31}=0$, consistent with $sI-A$ and all later computations ($A-\lambda_iI$, null spaces, the check of $AP$).\\
60 & \ref{ss:cht} & $A-4+3A^{-1}=0$ & $A-4I+3A^{-1}=0$ (identity matrix missing).\\
61 & \ref{ss:minpoly} & $m(A)=(A-2)(A-2)$ & $m(A)=(A-2I)(A-I)$.\\
61 & \ref{ss:minpoly} & $d(s)=m(s)d(s)$ & $d(s)=m(s)q(s)$.\\
61 & \ref{ss:minpoly} & Scattered list of crossed-out candidates with check marks (also on crossed-out ones and on $m(s)=s-1$); last candidate $(s-1)(s-2)^1$ & Organized by properties II and III: the only possibilities are $(s-1)(s-2)$ and $(s-1)(s-2)^2$.\\
61 & \ref{ss:minpoly} & Upper bound $m_i\le k$ & $m_i\le r_i$ (cf.\ p.~63).\\
62 & \ref{ss:minpoly} & $\bar r(s)=r(\alpha)/\alpha$; division $d=mq+r$ only implicit & $\bar r(s)=r(s)/\alpha$; the division $d(s)=m(s)q(s)+r(s)$ written out.\\
62 & \ref{ss:minpoly} & $e_1$ instead of $e_i$ (twice); stray cancelled $e_i$ inside the bracket & $e_i$; stray term removed.\\
63 & \ref{ss:minpoly-decomp} & Exponent 2 on $(s-\lambda_2)$ & All exponents are $1$ in the distinct-eigenvalue case.\\
63 & \ref{ss:minpoly-decomp} & $(A_2-I)(A_2-2I)=0$ & The product is nonzero, which is exactly why $m_2(s)=(s-1)(s-2)^2$.\\
%
\grp{sec:jordan}
65 & \ref{ss:gen-eigenspaces} & $\Null\big((A-2I)\big)$ in the last column & $\Null\big((A-2I)^3\big)$.\\
66 & \ref{ss:nullspace-chains} & ``$=$'' missing between the two null spaces & $\Null(\Sigma_i^l)=\Null(\Sigma_i^{l+1})$.\\
66 & \ref{ss:nullspace-chains} & $A$ instead of $A_1$ in the dimension column and last row & $A_1$.\\
67 & \ref{ss:counting-jb} & $\dim\Null(\Sigma_i^{l+1})-\dim\Null(\Sigma_i^{l})=\#$ JB's with size $\ge l$ (off by one) & $\dim\Null(\Sigma_i^{l})-\dim\Null(\Sigma_i^{l-1})=\#$ JB's with size $\ge l$; indices shifted so that $l=1$ is the first line.\\
69 & \ref{ss:chains} & $\Sigma_i^{m_i-1}\neq0$ (vector missing) & $\Sigma_i^{m_i-1}x\neq0$.\\
71 & \ref{ss:jordan-examples} & ``\# of JB w.s.\ $=2$ for $\lambda_2$'', ``\# of JB with size $=3$ for $\lambda_2$'' & These counts are for $\lambda_1$ and mean size $\ge2$, $\ge3$.\\
72 & \ref{ss:jordan-special} & ``$1\le1$'' under $\Null(\Sigma_i)$ & $\dim\Null(\Sigma_i)=1$ (each strict step gains exactly $+1$).\\
73 & \ref{ss:jordan-special} & $(\Sigma_i^{m_i-2}x)\Sigma_i$ (factor order); $(\lambda_i+\Sigma_i)$ & $\Sigma_i(\Sigma_i^{m_i-2}x)$; $(\lambda_iI+\Sigma_i)$.\\
%
\grp{sec:hermitian}
78 & \ref{ss:hermitian-props} & $\lambda_i=\lambda_j\Rightarrow\inner{e_i}{e_j}=0$ & $\lambda_i\neq\lambda_j$: the hypothesis is distinct eigenvalues.\\
79 & \ref{ss:hermitian-props} & $1\le m_i\le r_k$ & $1\le m_i\le r_i$.\\
80 & \ref{ss:hermitian-props} & ``For diagonalization, we need'' normality; $N^2=N^2$ & Normality is needed for \emph{unitary} diagonalization, not diagonalization in general; $N\Hm N=N^2=NN\Hm$ when $N\Hm=N$.\\
80 & \ref{ss:hermitian-props} & ``$\Leftrightarrow$ $A$ is normal'' & ``$\Longleftarrow$'': normality implies the three statements, but a diagonalizable matrix need not be normal.\\
83 & \ref{ss:posdef} & ``$\norm{e_i}^2$ is non-negative'' & Positive (as $e_i\neq0$): strict positivity is what is needed for $\lambda_i>0$.\\
%
\grp{sec:functions}
84 & \ref{ss:fn-def} & $\frac{d^if}{dt^i}$ & $\frac{d^if}{ds^i}$: the variable is $s$.\\
86 & \ref{ss:method1-diag} & ``Eigenvalues are invariant under functions'' & Eigenvectors $e_j$ are invariant; the eigenvalues become $f(\lambda_j)$.\\
86 & \ref{ss:method1-diag} & $\log 0$ & $\log 1=0$: the eigenvalue is $1$.\\
86 & \ref{ss:method1-diag} & $c_2=-\frac{\log3}{2}$ & $c_0=-\frac{\log3}{2}$.\\
86 & \ref{ss:method1-diag} & Left-hand side $\log s$ & $\log A$.\\
86 & \ref{ss:method1-diag} & $f'(\lambda_j)=p(\lambda_j)$ (as written) & $f(\lambda_j)=p(\lambda_j)$.\\
86 & \ref{ss:method1-diag} & $\dot x=\bmat{2&1\\1&2}x_0$ & $\dot x=Ax$ with $x(0)=x_0$, so $x(t)=e^{At}x_0$.\\
89 & \ref{ss:method1-general} & (Earlier transcription) $c_1$ struck through in $p(1)$ & $c_1=\pi\neq0$ is substituted, giving $c_3+c_4=-\pi$.\\
89 & \ref{ss:method1-general} & $f^{(t)}(A_j)$ & $f^{(t)}(\lambda_j)$.\\
89 & \ref{ss:method1-general} & ``Square root function DNA.'' & Read as ``does not apply'': $f'(0)$ does not exist (consistent with the remark above).\\
90 & \ref{ss:method2} & Generic block labelled $J_\sigma$ (size $d_t$) and last block of $J^i$ labelled $J_t^i$ & Made consistent: blocks $J_1,\dots,J_\sigma$, generic block $J_t$.\\
92 & \ref{ss:method2} & $\mp$ in both diagonal entries of $A^{1/2}$ & $\pm\sqrt3$, $\pm\sqrt1$: the two signs are chosen independently (4 square roots).\\
\end{longtable}
}
